% ============================================================
%  Tất cả công thức toán học từ REPORT_FULL_VI.md
%  Chuyển đổi sang LaTeX
% ============================================================

\documentclass[12pt, a4paper]{article}
\usepackage[utf8]{inputenc}
\usepackage[vietnamese]{babel}
\usepackage{amsmath, amssymb, amsthm}
\usepackage{geometry}
\geometry{margin=2.5cm}

\title{Công thức toán học --- Đề tài Mini-Project 09:\\
2D Image Processing \& Video Motion Analysis}
\author{DSP --- Đại học Bách khoa Hà Nội}
\date{}

\begin{document}
\maketitle
\tableofcontents
\newpage

% ============================================================
\section{Biến đổi Fourier rời rạc 2D (DFT 2D)}

\subsection{Định nghĩa DFT 2D thuận}

Cho ảnh rời rạc $f(x,y)$ kích thước $M \times N$, với $x \in \{0,\ldots,M-1\}$,
$y \in \{0,\ldots,N-1\}$:

\begin{equation}
  F(u,v) = \sum_{x=0}^{M-1} \sum_{y=0}^{N-1}
            f(x,y)\, W_M^{\,ux}\, W_N^{\,vy}
  \label{eq:dft2d}
\end{equation}

trong đó các \emph{twiddle factors} được định nghĩa là:
\begin{equation}
  W_M = e^{-j2\pi/M}, \qquad W_N = e^{-j2\pi/N}
  \label{eq:twiddle}
\end{equation}

và $u \in \{0,\ldots,M-1\}$, $v \in \{0,\ldots,N-1\}$ là các chỉ số
tần số không gian.

\subsection{DFT 2D nghịch đảo (IDFT 2D)}

\begin{equation}
  f(x,y) = \frac{1}{MN}
            \sum_{u=0}^{M-1} \sum_{v=0}^{N-1}
            F(u,v)\, W_M^{-ux}\, W_N^{-vy}
  \label{eq:idft2d}
\end{equation}

Thành phần DC: $F(0,0) = \displaystyle\sum_x \sum_y f(x,y)$.

\subsection{Tính phân tách của nhân DFT 2D}

Nhân DFT 2D phân tách thành tích của hai nhân DFT 1D độc lập:
\begin{equation}
  W_M^{\,ux} \cdot W_N^{\,vy} = e^{-j2\pi ux/M} \cdot e^{-j2\pi vy/N}
  \label{eq:sep}
\end{equation}

\subsection{Phân rã hàng--cột}

\textbf{Bước 1} --- DFT 1D theo từng hàng:
\begin{equation}
  G(x,v) = \sum_{y=0}^{N-1} f(x,y)\, W_N^{\,vy}, \quad \forall\, x
  \label{eq:row_dft}
\end{equation}

\textbf{Bước 2} --- DFT 1D theo từng cột của $G$:
\begin{equation}
  F(u,v) = \sum_{x=0}^{M-1} G(x,v)\, W_M^{\,ux}, \quad \forall\, v
  \label{eq:col_dft}
\end{equation}

Độ phức tạp thời gian toàn bộ:
$\mathcal{O}(MN\log N) + \mathcal{O}(NM\log M) = \mathcal{O}(MN\log(MN))$.

% ============================================================
\section{Các bộ lọc trong miền tần số}

\subsection{Phép lọc tổng quát}

Đầu ra đã lọc trong miền không gian:
\begin{equation}
  g(x,y) = \mathcal{F}^{-1}\!\bigl[F(u,v) \cdot H(u,v)\bigr]
  \label{eq:filter_general}
\end{equation}

Tương đương tích chập 2D $g = h * f$ (Định lý tích chập), nhưng tính
trong $\mathcal{O}(MN\log MN)$ thay vì $\mathcal{O}(M^2N^2)$.

\subsection{Hàm khoảng cách bán kính}

Khoảng cách từ điểm $(u,v)$ đến tâm phổ $(M/2,\, N/2)$:
\begin{equation}
  D(u,v) = \sqrt{\!\left(u - \tfrac{M}{2}\right)^{\!2} + \left(v - \tfrac{N}{2}\right)^{\!2}}
  \label{eq:dist}
\end{equation}

\subsection{Bộ lọc thông thấp lý tưởng (ILP)}

\begin{equation}
  H_{\mathrm{ILP}}(u,v) = \begin{cases} 1 & \text{nếu } D(u,v) \le D_0  0 & \text{nếu } D(u,v) > D_0 \end{cases}
  \label{eq:ilp}
\end{equation}

\subsection{Bộ lọc thông thấp Butterworth (bậc $n$)}

\begin{equation}
  H_{\mathrm{BW}}(u,v) = \frac{1}{1 + \left(\dfrac{D(u,v)}{D_0}\right)^{2n}}
  \label{eq:bw}
\end{equation}

Tại $D = D_0$: $H_{\mathrm{BW}} = 0{,}5$ (điểm $-3\,\mathrm{dB}$),
bất kể bậc $n$.

\subsection{Bộ lọc thông cao lý tưởng (IHP)}

\begin{equation}
  H_{\mathrm{IHP}}(u,v) = 1 - H_{\mathrm{ILP}}(u,v)
  \label{eq:ihp}
\end{equation}

\subsection{Bộ lọc thông cao Gauss}

\begin{equation}
  H_{\mathrm{GHP}}(u,v) = 1 - \exp\!\left(-\frac{D(u,v)^2}{2D_0^2}\right)
  \label{eq:ghp}
\end{equation}

Tại $D = D_0$: $H_{\mathrm{GHP}} = 1 - e^{-1/2} \approx 0{,}394$
(điểm $\approx -4\,\mathrm{dB}$).

\subsection{Bộ lọc khấc (Notch filter)}

\begin{equation}
  H_{\mathrm{notch}}(u,v) =
  \begin{cases}
    0 & \text{nếu } D\!\left((u,v),\, (\pm u_0, \pm v_0)\right) \le r \\
    1 & \text{nếu khác}
  \end{cases}
  \label{eq:notch}
\end{equation}

Các cặp liên hợp $(\pm u_0, \pm v_0)$ đặt đối xứng để bảo toàn
$F^*(u,v) = F(-u,-v)$, đảm bảo đầu ra có giá trị thực.

% ============================================================
\section{Mô hình suy giảm chất lượng ảnh}

\subsection{Khung suy giảm--phục hồi}

Trong miền không gian:
\begin{equation}
  g(x,y) = (h * f)(x,y) + n(x,y)
  \label{eq:degradation_spatial}
\end{equation}

Trong miền tần số (Định lý tích chập):
\begin{equation}
  G(u,v) = H(u,v) \cdot F(u,v) + N(u,v)
  \label{eq:degradation_freq}
\end{equation}

$H(u,v)$ là Hàm truyền đạt quang học (OTF) --- biến đổi Fourier của
Hàm lan rộng điểm (PSF) $h(x,y)$.

\subsection{PSF mờ chuyển động}

Chuyển động ngang $L$ điểm ảnh ở tốc độ đều:
\begin{equation}
  h_{\mathrm{motion}}(x,y) = 
  \begin{cases}
    \dfrac{1}{L} & \text{nếu } y = 0 \text{ và } 0 \le x < L \\[6pt]
    0            & \text{nếu khác}
  \end{cases}
  \label{eq:psf_motion}
\end{equation}

OTF tương ứng:
\begin{equation}
  H_{\mathrm{motion}}(u,v) = \frac{1}{L}\sum_{k=0}^{L-1} e^{-j2\pi ku/M} = \operatorname{sinc}\!\left(\frac{uL}{M}\right) \cdot e^{-j\pi u(L-1)/M}
  \label{eq:otf_motion}
\end{equation}

Các \emph{điểm không} (zeros) của $H_{\mathrm{motion}}$ xuất hiện tại
$u = kM/L$ với $k \in \mathbb{Z}$, $k \ne 0$.

\subsection{PSF mờ mất tiêu điểm (đĩa tròn bán kính $r$)}

\begin{equation}
  h_{\mathrm{defocus}}(x,y) =
  \begin{cases}
    \dfrac{1}{\pi r^2} & \text{nếu } x^2 + y^2 \le r^2 \\[6pt]
    0                  & \text{nếu khác}
  \end{cases}
  \label{eq:psf_defocus}
\end{equation}

% ============================================================
\section{Giải tích chập Wiener}

\subsection{Ước lượng MMSE}

Bộ lọc Wiener là \emph{ước lượng tuyến tính sai số bình phương trung
bình tối thiểu} (MMSE) của $f$ từ $g$:
\begin{equation}
  \hat{F}(u,v) = W(u,v) \cdot G(u,v)
  \label{eq:wiener_out}
\end{equation}

Hàm truyền đạt Wiener tối ưu:
\begin{equation}
  W(u,v)
  = \frac{H^*(u,v)\, S_{ff}(u,v)}
         {\lvert H(u,v)\rvert^2 S_{ff}(u,v) + S_{nn}(u,v)}
  = \frac{H^*(u,v)}
         {\lvert H(u,v)\rvert^2 + S_{nn}(u,v)/S_{ff}(u,v)}
  \label{eq:wiener_full}
\end{equation}

trong đó $S_{ff}$ và $S_{nn}$ lần lượt là mật độ phổ công suất của
tín hiệu và nhiễu.

\subsection{Xấp xỉ hằng số $K$}

Thay thế tỷ số $S_{nn}/S_{ff}$ bằng hằng số vô hướng $K$:
\begin{equation}
  W(u,v) = \frac{H^*(u,v)}{\lvert H(u,v)\rvert^2 + K}
  \label{eq:wiener_K}
\end{equation}

Giới hạn đặc biệt:
\begin{itemize}
  \item $K \to 0$: bộ lọc nghịch đảo thuần túy
        $W \approx 1/H$ --- khuếch đại vô hạn tại các điểm không OTF.
  \item $K \to \infty$: $W \to 0$ --- không phục hồi gì.
  \item $K = \sigma_n^2 / \sigma_f^2$: tối ưu MSE cho nhiễu trắng
        và phổ tín hiệu phẳng.
\end{itemize}

Hành vi tại hai vùng tần số:
\begin{align}
  \lvert H\rvert^2 \gg K &\implies W(u,v) \approx \frac{1}{H(u,v)}
    \quad \text{(nghịch đảo hoàn hảo)} \label{eq:wiener_high_snr} \\
  \lvert H\rvert^2 \ll K &\implies W(u,v) \to 0
    \quad \text{(triệt nhiễu)} \label{eq:wiener_low_snr}
\end{align}

\subsection{Lựa chọn $K$ theo hướng dẫn SSIM}

\begin{equation}
  K^* = \operatorname*{arg\,max}_{K \,\in\, \mathcal{K}}
        \operatorname{SSIM}\!\bigl(\hat{f}(K),\, f_{\mathrm{gốc}}\bigr)
  \label{eq:K_optimal}
\end{equation}

với tập ứng viên $\mathcal{K} = \{10^{-5},\,10^{-4},\,10^{-3},\,10^{-2},\,10^{-1}\}$.

\subsection{Chỉ số tương đồng cấu trúc (SSIM)}

\begin{equation}
  \operatorname{SSIM}(x,y)
  = \bigl[l(x,y)\bigr]^{\alpha}
    \bigl[c(x,y)\bigr]^{\beta}
    \bigl[s(x,y)\bigr]^{\gamma}
  \label{eq:ssim}
\end{equation}

với ba thành phần:
\begin{align}
  l(x,y) &= \frac{2\mu_x\mu_y + C_1}{\mu_x^2 + \mu_y^2 + C_1}
             \quad \text{(so sánh độ sáng)} \label{eq:ssim_l} \\[4pt]
  c(x,y) &= \frac{2\sigma_x\sigma_y + C_2}{\sigma_x^2 + \sigma_y^2 + C_2}
             \quad \text{(so sánh độ tương phản)} \label{eq:ssim_c} \\[4pt]
  s(x,y) &= \frac{\sigma_{xy} + C_3}{\sigma_x\sigma_y + C_3}
             \quad \text{(so sánh cấu trúc)} \label{eq:ssim_s}
\end{align}

Hằng số ổn định ($L = 1{,}0$ cho ảnh chuẩn hóa):
\begin{equation}
  C_1 = (0{,}01 L)^2, \quad
  C_2 = (0{,}03 L)^2, \quad
  C_3 = \frac{C_2}{2}
  \label{eq:ssim_const}
\end{equation}

Với $\alpha = \beta = \gamma = 1$: $\operatorname{SSIM} \in [-1,\,1]$,
bằng $1$ khi hai ảnh giống hệt nhau.

% ============================================================
\section{Trích xuất đặc trưng kết cấu và phân loại KNN}

\subsection{Phổ công suất 2D}

\begin{equation}
  P(u,v) = \lvert F(u,v) \rvert^2
  \label{eq:power_spectrum}
\end{equation}

Để hiển thị, phổ công suất log được dùng:
$\log\!\bigl(1 + P(u,v)\bigr)$.

\subsection{Đặc trưng năng lượng theo hướng}

Tích phân công suất qua ô góc $[\theta,\, \theta + \Delta\theta)$
(loại trừ DC, $r > 2$):
\begin{equation}
  E_{\mathrm{dir}}(\theta)
  = \sum_{\substack{(u,v)\,:\\ \angle(u,v)\,\in\,[\theta,\,\theta+\Delta\theta)}}
    P(u,v)
  \label{eq:feat_dir}
\end{equation}

$36$ ô, chiều rộng $\Delta\theta = 5°$ $\Rightarrow$ vectơ đặc trưng $36$ chiều.

\subsection{Đặc trưng năng lượng theo bán kính}

Tích phân công suất qua vành đồng tâm $[r,\, r + \Delta r)$:
\begin{equation}
  E_{\mathrm{rad}}(r)
  = \sum_{\substack{(u,v)\,:\\ D(u,v)\,\in\,[r,\,r+\Delta r)}}
    P(u,v)
  \label{eq:feat_rad}
\end{equation}

$32$ ô bán kính $\Rightarrow$ vectơ đặc trưng $32$ chiều.

\textbf{Vectơ đặc trưng kết hợp:} $36 + 32 = 68$ chiều,
được chuẩn hóa $L_2$.

\subsection{Bộ phân loại $K$ lân cận gần nhất (KNN)}

Cho tập huấn luyện $\{(\mathbf{x}_i, y_i)\}$ và mẫu thử $\mathbf{x}^*$:
\begin{equation}
  y^* = \operatorname{majority\_vote}
        \bigl\{\, y_i : i \in \mathcal{N}_k(\mathbf{x}^*) \,\bigr\}
  \label{eq:knn}
\end{equation}

$\mathcal{N}_k(\mathbf{x}^*)$ là tập $k$ lân cận gần nhất theo khoảng
cách Euclid trong không gian đặc trưng. Đánh giá bằng kiểm định chéo
bỏ một (LOO-CV).

% ============================================================
\section{Luồng quang học Horn--Schunck}

\subsection{Phát biểu bài toán}

Cho hai khung hình xám liên tiếp $I_1(x,y)$ và $I_2(x,y)$, tìm trường
luồng $(u(x,y),\, v(x,y))$ sao cho:
\begin{equation}
  I_1(x,y) \approx I_2\!\bigl(x + u(x,y),\; y + v(x,y)\bigr)
  \label{eq:of_problem}
\end{equation}

\subsection{Ràng buộc bất biến độ sáng}

Khai triển Taylor bậc nhất:
\begin{equation}
  I_2(x+u,\, y+v) \approx I_1(x,y) + I_x u + I_y v + I_t
  \label{eq:taylor}
\end{equation}

\textbf{Phương trình ràng buộc luồng quang học (OFCE)}:
\begin{equation}
  I_x u + I_y v + I_t = 0
  \label{eq:ofce}
\end{equation}

trong đó:
\begin{align*}
  I_x &= \partial I/\partial x \quad \text{(gradient không gian ngang)} \\
  I_y &= \partial I/\partial y \quad \text{(gradient không gian dọc)} \\
  I_t &= I_2 - I_1 \quad \text{(gradient thời gian)}
\end{align*}

\subsection{Chính quy hóa trơn Horn--Schunck}

Horn và Schunck (1981) tối thiểu hóa phiến năng lượng kết hợp:
\begin{equation}
  E(u,v)
  = \iint \left[
      \bigl(I_x u + I_y v + I_t\bigr)^2
      + \alpha^2\!\left(\lvert\nabla u\rvert^2 + \lvert\nabla v\rvert^2\right)
    \right] dx\, dy
  \label{eq:hs_energy}
\end{equation}

\begin{itemize}
  \item Số hạng thứ nhất: độ trung thực dữ liệu (phạt vi phạm OFCE).
  \item Số hạng thứ hai: chính quy hóa trơn (phạt biến thiên luồng).
  \item $\alpha > 0$: trọng số trơn.
\end{itemize}

\subsection{Các phương trình Euler--Lagrange}

Đặt đạo hàm biến phân bằng không:
\begin{align}
  \alpha^2 \Delta u &= I_x\!\left(I_x u + I_y v + I_t\right) \label{eq:el_u} \\
  \alpha^2 \Delta v &= I_y\!\left(I_x u + I_y v + I_t\right) \label{eq:el_v}
\end{align}

trong đó toán tử Laplacian rời rạc xấp xỉ:
$\Delta u \approx \bar{u} - u$ ($\bar{u}$: trung bình không gian cục bộ).

\subsection{Nghiệm lặp (Gauss--Seidel)}

\begin{align}
  u^{n+1} &= \bar{u}^n
    - \frac{I_x\!\left(I_x \bar{u}^n + I_y \bar{v}^n + I_t\right)}
           {\alpha^2 + I_x^2 + I_y^2}
  \label{eq:hs_iter_u} \\[6pt]
  v^{n+1} &= \bar{v}^n
    - \frac{I_y\!\left(I_x \bar{u}^n + I_y \bar{v}^n + I_t\right)}
           {\alpha^2 + I_x^2 + I_y^2}
  \label{eq:hs_iter_v}
\end{align}

\subsection{Nhân lấy trung bình Horn--Schunck}

$\bar{u}$ và $\bar{v}$ được tính bằng tích chập với nhân $3\times 3$:
\begin{equation}
  K_{\mathrm{HS}} =
  \begin{bmatrix}
    \tfrac{1}{12} & \tfrac{1}{6} & \tfrac{1}{12} \\[4pt]
    \tfrac{1}{6}  & 0            & \tfrac{1}{6}  \\[4pt]
    \tfrac{1}{12} & \tfrac{1}{6} & \tfrac{1}{12}
  \end{bmatrix}
  \label{eq:hs_kernel}
\end{equation}

(Trọng số tâm bằng $0$ --- trung bình $8$ lân cận có trọng số.)

\subsection{Tính gradient tại khung trung bình}

\begin{align}
  I_{\mathrm{avg}} &= \frac{I_1 + I_2}{2} \label{eq:grad_avg} \\[4pt]
  I_x &= \frac{\partial I_{\mathrm{avg}}}{\partial x}
         \quad \text{(sai phân trung tâm)} \label{eq:grad_x} \\[4pt]
  I_y &= \frac{\partial I_{\mathrm{avg}}}{\partial y} \label{eq:grad_y} \\[4pt]
  I_t &= I_2 - I_1 \label{eq:grad_t}
\end{align}

\subsection{Biên độ trường luồng}

\begin{equation}
  M(x,y) = \sqrt{u(x,y)^2 + v(x,y)^2}
  \label{eq:flow_magnitude}
\end{equation}

% ============================================================
\section{Đường dẫn phân tích chuyển động video}

Biên độ luồng trung bình tại cặp khung $i$:
\begin{equation}
  m_i = \frac{1}{MN}\sum_{x,y} \sqrt{u_i(x,y)^2 + v_i(x,y)^2}
  \label{eq:mean_motion}
\end{equation}

Cặp khung đỉnh:
\begin{equation}
  i^* = \operatorname*{arg\,max}_{i}\; m_i
  \label{eq:peak_frame}
\end{equation}

% ============================================================
\section*{Tổng hợp ký hiệu}

\begin{center}
\renewcommand{\arraystretch}{1.3}
\begin{tabular}{ll}
\hline
\textbf{Ký hiệu} & \textbf{Ý nghĩa} \\
\hline
$f(x,y)$          & Ảnh gốc trong miền không gian \\
$F(u,v)$          & DFT 2D của $f$ \\
$g(x,y)$          & Ảnh bị suy giảm (sau mờ + nhiễu) \\
$h(x,y)$, $H(u,v)$& PSF và OTF \\
$n(x,y)$, $N(u,v)$& Nhiễu trong miền không gian / tần số \\
$W(u,v)$          & Hàm truyền đạt bộ lọc Wiener \\
$D(u,v)$          & Khoảng cách bán kính từ DC \\
$D_0$             & Tần số cắt \\
$K$               & Tham số chính quy hóa Wiener \\
$S_{ff}$, $S_{nn}$& Mật độ phổ công suất tín hiệu / nhiễu \\
$P(u,v)$          & Phổ công suất $|F(u,v)|^2$ \\
$I_x$, $I_y$, $I_t$& Gradient không gian và thời gian \\
$u(x,y)$, $v(x,y)$& Thành phần trường luồng quang học \\
$\alpha$          & Trọng số trơn Horn--Schunck \\
$\bar{u}$, $\bar{v}$& Trung bình không gian cục bộ của luồng \\
\hline
\end{tabular}
\end{center}

\end{document}
